% GENERATED FILE. Edit canonical lessons, then run npm run generate:books.
% canonical-source-hash: fnv1a64:98323d000c1e590d
% canonical-lessons: ZH-C118-baoguo, ZH-C118-ji, ZH-C118-xin, ZH-C118-duanxin, ZH-C118-shou

\chapter{Sending Things}
\label{ch:zh-sending-things}
\hlchaptermodality{118}

\begin{chapteropening}
\textbf{By the end of this chapter:} \emph{I can post a parcel, talk about letters and text messages, and say receive.}

Complete the last lesson of chapter 118: I can post a parcel, talk about letters and text messages, and say receive.
\end{chapteropening}

\section[bāoguǒ]{bāoguǒ — parcel; package}

\label{lesson:ZH-C118-baoguo}



\begin{quote}
Before the new one: say the Mandarin for the north, then the Mandarin for direction.
\end{quote}

\subsection*{You'll want to know: bāoguǒ}
\emph{bāoguǒ} — \textquotedblleft{}parcel; package\textquotedblright{}.

\begin{quote}
\emph{zhège bāoguǒ shì nǐ de ma} — Is this parcel yours?
\end{quote}

Its characters come later; for now it is a word for the ear.

\begin{cousinweb}
\emph{bāoguǒ} is built from bāo and guǒ, both inherited from Middle Chinese, \emph{paew} and \emph{kwaX} (Guangyun, in Baxter's transcription); the dictionary glosses the second 'to wrap, to bind', and gives the pair itself as the gloss of the first.
\end{cousinweb}

\begin{grammarlens}[title={What you've built}]
One more word for everyday actions.
\end{grammarlens}

\subsection*{Your turn}
\begin{itemize}
\raggedright
  \item \emph{Say it:} \emph{bāoguǒ}
  \item \emph{Say it:} \emph{bāoguǒ}, once more
  \item \emph{Say it:} \emph{zhège bāoguǒ shì nǐ de ma} — Is this parcel yours?
  \item \emph{Recall:} say the Mandarin for the north, then the Mandarin for direction, then say \emph{bāoguǒ} again
\end{itemize}

\begin{tcolorbox}[breakable,colback=teal!4,colframe=teal!35!black,title={Before you move on}]
What does bāoguǒ mean? (\textquotedblleft{}Parcel; package\textquotedblright{}.) Say it once more.
\end{tcolorbox}
\section[jì]{jì — to post; to send by post}

\label{lesson:ZH-C118-ji}



\begin{quote}
Before the new one: say the Mandarin for direction, then the Mandarin for parcel.
\end{quote}

\subsection*{You'll want to know: jì}
\emph{jì} — \textquotedblleft{}to post; to send by post\textquotedblright{}.

\begin{quote}
\emph{wǒ qù yóujú jì bāoguǒ} — I'm going to the post office to send a parcel.
\end{quote}

Its characters come later; for now it is a word for the ear.

\begin{cousinweb}
\emph{jì} is inherited from Middle Chinese \emph{kjeH} (Guangyun, in Baxter's transcription), glossed 'to entrust to someone'.
\end{cousinweb}

\begin{grammarlens}[title={What you've built}]
One more word for everyday actions.
\end{grammarlens}

\subsection*{Your turn}
\begin{itemize}
\raggedright
  \item \emph{Say it:} \emph{jì}
  \item \emph{Say it:} \emph{jì}, once more
  \item \emph{Say it:} \emph{wǒ qù yóujú jì bāoguǒ} — I'm going to the post office to send a parcel.
  \item \emph{Recall:} say the Mandarin for direction, then the Mandarin for parcel, then say \emph{jì} again
\end{itemize}

\begin{tcolorbox}[breakable,colback=teal!4,colframe=teal!35!black,title={Before you move on}]
What does jì mean? (\textquotedblleft{}To post; to send by post\textquotedblright{}.) Say it once more.
\end{tcolorbox}
\section[xìn]{xìn — letter (in the post)}

\label{lesson:ZH-C118-xin}



\begin{quote}
Before the new one: say the Mandarin for parcel, then the Mandarin for to post.
\end{quote}

\subsection*{You'll want to know: xìn}
\emph{xìn} — \textquotedblleft{}letter (in the post)\textquotedblright{}.

\begin{quote}
\emph{wǒ xiǎng gěi péngyou xiě xìn} — I'd like to write to a friend.
\end{quote}

Its characters come later; for now it is a word for the ear.

\begin{cousinweb}
\emph{xìn} is inherited from Middle Chinese \emph{sinH} (Guangyun, in Baxter's transcription), glossed 'trust' and 'proof'; how it came to mean a letter is not settled.
\end{cousinweb}

\begin{grammarlens}[title={What you've built}]
One more word for everyday actions.
\end{grammarlens}

\subsection*{Your turn}
\begin{itemize}
\raggedright
  \item \emph{Say it:} \emph{xìn}
  \item \emph{Say it:} \emph{xìn}, once more
  \item \emph{Say it:} \emph{wǒ xiǎng gěi péngyou xiě xìn} — I'd like to write to a friend.
  \item \emph{Recall:} say the Mandarin for parcel, then the Mandarin for to post, then say \emph{xìn} again
\end{itemize}

\begin{tcolorbox}[breakable,colback=teal!4,colframe=teal!35!black,title={Before you move on}]
What does xìn mean? (\textquotedblleft{}Letter (in the post)\textquotedblright{}.) Say it once more.
\end{tcolorbox}
\section[duǎnxìn]{duǎnxìn — text message}

\label{lesson:ZH-C118-duanxin}



\begin{quote}
Before the new one: say the Mandarin for to post, then the Mandarin for letter.
\end{quote}

\subsection*{You'll want to know: duǎnxìn}
\emph{duǎnxìn} — \textquotedblleft{}text message\textquotedblright{}.

\begin{quote}
\emph{wǒ gěi māma xiě duǎnxìn} — I'm writing a text to my mum.
\end{quote}

Its characters come later; for now it is a word for the ear.

\begin{cousinweb}
\emph{duǎnxìn} is built from duǎn 'short' and xìn 'letter', both inherited from Middle Chinese, \emph{twanX} and \emph{sinH} (Guangyun, in Baxter's transcription): a short letter.
\end{cousinweb}

\begin{grammarlens}[title={What you've built}]
One more word for everyday actions.
\end{grammarlens}

\subsection*{Your turn}
\begin{itemize}
\raggedright
  \item \emph{Say it:} \emph{duǎnxìn}
  \item \emph{Say it:} \emph{duǎnxìn}, once more
  \item \emph{Say it:} \emph{wǒ gěi māma xiě duǎnxìn} — I'm writing a text to my mum.
  \item \emph{Recall:} say the Mandarin for to post, then the Mandarin for letter, then say \emph{duǎnxìn} again
\end{itemize}

\begin{tcolorbox}[breakable,colback=teal!4,colframe=teal!35!black,title={Before you move on}]
What does duǎnxìn mean? (\textquotedblleft{}Text message\textquotedblright{}.) Say it once more.
\end{tcolorbox}
\section[shōu]{shōu — to receive; to accept}

\label{lesson:ZH-C118-shou}



\begin{quote}
Before the new one: say the Mandarin for letter, then the Mandarin for text message.
\end{quote}

\subsection*{You'll want to know: shōu}
\emph{shōu} — \textquotedblleft{}to receive; to accept\textquotedblright{}.

\begin{quote}
\emph{tā bù shōu qián} — He won't take any money.
\end{quote}

Its characters come later; for now it is a word for the ear.

\begin{cousinweb}
\emph{shōu} is inherited from Middle Chinese \emph{syuw} (Guangyun, in Baxter's transcription), glossed 'to gather in' and 'to seize'.
\end{cousinweb}

\begin{grammarlens}[title={What you've built}]
One more word for everyday actions.
\end{grammarlens}

\subsection*{Your turn}
\begin{itemize}
\raggedright
  \item \emph{Say it:} \emph{shōu}
  \item \emph{Say it:} \emph{shōu}, once more
  \item \emph{Say it:} \emph{tā bù shōu qián} — He won't take any money.
  \item \emph{Recall:} say the Mandarin for letter, then the Mandarin for text message, then say \emph{shōu} again
\end{itemize}

\begin{tcolorbox}[breakable,colback=teal!4,colframe=teal!35!black,title={Before you move on}]
What does shōu mean? (\textquotedblleft{}To receive; to accept\textquotedblright{}.) Say it once more.
\end{tcolorbox}
